\section{Related work}
\label{sec:related}

\paragraph{Subtypes from co-occurring conditions.}
\citet{doshivelez2014} clustered diagnosis-code histories of about 4,900 patients at one hospital and found three small distinct groups (seizure, multisystem, psychiatric) and one large unresolved group, with a replication sample from a second region. \citet{lingren2016} reported psychiatric, developmental and seizure groupings in 20,658 hospital records. \citet{vargason2019} applied $k$-means to insurance claims for 3,278 young children and found high rates of everything (24\%), developmental delay (27\%) and low rates (50\%). \citet{montgomery2023} found four groups ordered by support need in 1,151 children from the Australian Autism Biobank, and \citet{litman2025} fit a generative mixture model to more than 5,000 SPARK participants and reported four classes (social and behavioral challenges 37\%, mixed with developmental delay 19\%, moderate challenges 34\%, broadly affected 10\%) that replicated in a second cohort and differed in genetic signal. Across these studies the recurring contrasts are a developmental or language group, a psychiatric or behavioral group and a small high-burden group; the number of groups varies with the inputs and the method.

\paragraph{Validation.}
\citet{agelink2021} distilled the ways a subtype can be validated into a seven-item checklist. External validation on variables not used for clustering was common (88\% of studies); independent replication (9\%), a second method (13\%) and parallel measures (6\%) were rare, and no study used all seven. Our checks follow the checklist (Section~\ref{sec:checks}). Stability under resampling is a separate tradition: \citet{benhur2002} and \citet{lange2004} choose the number of clusters by how well clusterings agree across subsamples, and \citet{hennig2007} gives thresholds for when a single cluster is stable. Current guidance on latent class analysis treats the number of classes as a judgment made from several indices together \citep{masyn2013,nylundgibson2018}.

\paragraph{Types or dimensions.}
\citet{bauer2003,bauer2004} showed that classes can stand in for continuous variation and that fit indices alone cannot tell the two apart. \citet{lubke2006} and \citet{markon2006} showed that fitting class and factor models to the same data on the same likelihood scale can, and \citet{lubke2005} proposed factor mixture models that allow both. Meta-analyses of taxometric studies in personality and psychopathology find dimensional structure about five times as often as categorical \citep{haslam2012,haslam2020}. For core autism traits, \citet{kim2019} found correlated dimensions fit caregiver ratings best in 3,825 clinic-referred youth, and \citet{wittkopf2023} found that BIC preferred dimensions for observation data but that factor scores were multimodal, and argued for a hybrid. We found no study that makes the comparison for co-occurring conditions. On the genetic side, \citet{zhang2025} found that common-variant autism genetics splits into two modestly correlated factors, one tied to earlier diagnosis and early communication difficulty and the other to later diagnosis, adolescent behavioral and emotional difficulty, ADHD and mental health. \citet{lombardo2019} argue that categorical and dimensional descriptions of autism heterogeneity need not exclude each other.

\paragraph{This study.}
The NSCH is a nationally representative survey with state-level samples \citep{ghandour2018}, used before to estimate the prevalence of parent-reported autism \citep{kogan2018} but not, as far as we found, to look for subtypes. Its class shares are population estimates rather than descriptions of a clinic sample. Its inputs are weaker than a clinic's: parent report of a provider's diagnosis, 18 yes/no items and one time point, which limits what it can say about any one child and rules out temporal and predictive validation.
